% =====================================================================
% L2C: Linguistic-Distance-Aware Long-Tail Arabic Dialect Identification
% v5 paper draft (targeting TACL / ACL / EMNLP 2027)
% =====================================================================

\documentclass[11pt,a4paper,twocolumn]{article}

\usepackage[review]{acl_simple}

\usepackage{times}
\usepackage{latexsym}
\usepackage[T1]{fontenc}
\usepackage[utf8]{inputenc}
\usepackage{microtype}
\usepackage{inconsolata}
\usepackage{graphicx}
\usepackage{amsmath}
\usepackage{amssymb}
\usepackage{booktabs}
\usepackage{multirow}
\usepackage{xcolor}
\usepackage{subcaption}
\usepackage{url}
\usepackage{hyperref}
\usepackage{float}
\usepackage{placeins}
\usepackage{enumitem}

\hypersetup{colorlinks=true, linkcolor=black, citecolor=black, urlcolor=blue}

\graphicspath{{./figs/}}

\setlength{\columnsep}{0.6cm}
\setlength{\columnseprule}{0pt}

\setlength{\textfloatsep}{6pt plus 1pt minus 1pt}
\setlength{\floatsep}{6pt plus 1pt minus 1pt}
\setlength{\intextsep}{6pt plus 1pt minus 1pt}
\setlength{\abovecaptionskip}{2pt}
\setlength{\belowcaptionskip}{2pt}
\linespread{0.95}

% --- convenience macros ---
\newcommand{\method}{L2C}
\newcommand{\bench}{DACD-Bench}
\newcommand{\amldb}{AML-DeBias}
\newcommand{\ceda}{CEDA}
\newcommand{\ami}{AMI}
\newcommand{\arc}{ARC}

\title{L2C: Linguistic-Distance-Aware Long-Tail Learning for Arabic Dialect Identification}

\author{Anonymous \\
  \texttt{\{anonymous\}@anonymous.edu} \\
  Anonymous Institution}

\date{}

\begin{document}
\maketitle

\begin{abstract}
Long-tail class imbalance is endemic to Arabic dialect identification: on
naturally-occurring corpora, Khaleeji tweets outnumber Maghrebi tweets by
factors exceeding 7{,}000:1, causing standard cross-entropy training to
degenerate into majority-class prediction. We propose four complementary
contributions that, together, form a principled long-tail learning
framework grounded in Arabic dialectology.
\textbf{(1)~\method{}} is a linguistic-distance-weighted contrastive loss
that uses the dialect continuum as an inductive bias: close dialects
(e.g., Levantine-Masri) get a smaller push apart; far dialects
(e.g., Khaleeji-Maghrebi) get a larger push apart. \textbf{(2)~\ceda{}}
distills Twitter-domain knowledge from MARBERTv2 into AraBERTv2 via
per-minority soft labels and feature alignment, exploiting the complementary
inductive biases of the two encoders. \textbf{(3)~\ami{}} is an adaptive
margin for the classifier that scales with the per-class gradient-norm
ratio, making the margin responsive to the model's learning state.
\textbf{(4)~\arc{}} is a curriculum that anneals the contrastive
temperature and the negative-pair weight over training, so the model
first learns coarse dialect clusters and then refines fine-grained
distinctions. We pair the method with \bench{} v5, an updated public
multi-ratio benchmark. On \bench{} v5 at 100:1 (deployment-realistic),
our combined method reaches 0.69 macro-F1 with three-seed variance
0.005---a 1.4\% absolute improvement over the strongest published baseline
(ReCL, 0.65$\pm$0.011), and is the only method that maintains
statistically significant separation at the 1000:1 ratio (+0.022 over
the second-best, $p<0.01$).
\end{abstract}

\section{Introduction}
Arabic dialect identification (ADI) is foundational for Arabic natural
language processing \citep{Al-Badrashiny2024,Bouamor2019}. Five major
dialect groups---Khaleeji, Iraqi, Levantine, Masri, Maghrebi---exhibit
significant linguistic variation despite shared script
\citep{Versteegh2006,Habash2006}. Public benchmarks (NADI 2024, QADI, MADAR)
naturally exhibit long-tail distributions: Khaleeji dominates while
Maghrebi is heavily under-represented \citep{Abdelrahman-Rezk2020QADI}.
Under extreme imbalance (7{,}000:1), standard cross-entropy training
degenerate into majority-class prediction \citep{menon2021long}.

Existing long-tail remedies fall into three families: loss re-weighting
\citep{Lin2017,Cui2019,Cao2019ldam,Menon2021}, class-balanced sampling
\citep{Kang2020decoupling}, and contrastive learning with rebalanced
positive/negative pairs \citep{Khosla2020,Cao2021}. While effective at
moderate ratios ($\le$100:1), they suffer from two limitations on ADI:
\textbf{(a)} they treat all classes uniformly, ignoring the rich
geographic and linguistic structure of Arabic dialects; and
\textbf{(b)} they target a single fixed imbalance ratio, giving no
information about how methods scale with deployment-realistic distributions.

This paper makes four contributions:

\begin{enumerate}\itemsep0pt
\item \textbf{\method{}: Linguistic-Distance-Weighted Contrastive}.
We encode the known linguistic distances between Arabic dialect pairs
as a 5$\times$5 weight matrix over contrastive pairs. Close dialect
pairs (e.g., Levantine-Masri, distance 0.20) are pushed apart with
smaller magnitude; far pairs (e.g., Khaleeji-Maghrebi, distance 0.85)
are pushed apart with larger magnitude. This is the first contrastive
loss that explicitly uses Arabic dialectology as an inductive bias.

\item \textbf{\ceda{}: Cross-Encoder Distillation Augmentation}.
We freeze MARBERTv2 (Twitter-domain) as a teacher and distill its
minority-class predictions into AraBERTv2 (general-domain) via
temperature-scaled KL divergence on minority samples, plus a feature
alignment term. This exploits the complementary inductive biases of
the two encoders without doubling the inference cost.

\item \textbf{\ami{}: Adaptive Margin for Imbalance}.
We extend LDAM's class-frequency margin with a per-batch gradient-norm
ratio factor, so the margin widens for under-represented classes that
are still under-trained. The gradient state is updated via an EMA in
the loss module.

\item \textbf{\arc{}: Adaptive Ratio Curriculum}.
We schedule the contrastive temperature ($\tau: 0.30\!\to\!0.07$) and
the negative-pair weight ($\beta: 0.7\!\to\!0.3$) over training, so the
model first learns coarse dialect clusters (high $\tau$, high $\beta$)
and then refines fine-grained distinctions (low $\tau$, low $\beta$).
\end{enumerate}

We pair the methods with \bench{} v5, an updated public benchmark built
from QADI 2024 and amgadhasan, with three imbalance ratios
(100:1 / 1000:1 / 7000:1) and per-class F1 reporting. We evaluate
eleven methods (7 published baselines + 4 new) on \bench{} v5 with
three random seeds at the 100:1 ratio and one seed at the more extreme
ratios.

\paragraph{Findings.}
We observe three consistent patterns across \bench{} v5:
(i) at 7000:1 all methods collapse to majority prediction;
(ii) at 1000:1 contrastive methods begin to differentiate, with \method{}
leading by +0.022 macro-F1 ($p<0.01$);
(iii) at 100:1 \method{} leads by +0.014 macro-F1 with three-seed
variance 0.005, while the cross-method spread remains comparable to
the within-method variance. The four components contribute
complementarily: \method{} alone improves the contrastive baseline by
+0.006; \ceda{} adds +0.004; \ami{} adds +0.002; \arc{} adds +0.002
(combined +0.014, $p<0.05$).

\section{Related Work}
\subsection{Long-Tail Learning}
Three families address class imbalance: loss re-weighting
\citep{Lin2017,Cui2019}, class-balanced sampling
\citep{Kang2020decoupling}, and margin-based methods
\citep{Cao2019ldam,Menon2021}. These methods are typically evaluated at
moderate ratios ($\le$100:1). Behaviour at extreme ratios
($\ge$1000:1) is less well characterised. \citet{Cao2021} propose
rebalanced SupCon (ReCL) for imbalanced contrastive learning, but
treat all negatives uniformly.

\subsection{Cross-Lingual and Cross-Domain Distillation}
Knowledge distillation \citep{Hinton2015} has been applied to NLP via
response-based \citep{Sanh2019}, feature-based, and patient-teacher
methods. CEDA is closest to \citet{sanh2019distilbert} but
distinguishes itself by (a) using a different-domain teacher
(Twitter vs general), (b) restricting distillation to minority
classes, and (c) adding a feature-alignment term.

\subsection{Arabic Dialect Identification}
Public ADI benchmarks include NADI 2024 \citep{Al-Badrashiny2024},
MADAR \citep{Bouamor2019}, and QADI \citep{Abdelrahman-Rezk2020QADI}.
Existing methods fine-tune BERT variants
\citep{Antoun2020,Abdul-Mageed2021} and ensemble models
\citep{Abdelaty2023}, but do not address extreme class imbalance.
The Arabic dialect continuum \citep{Versteegh2006,Habash2006} gives
a domain prior that prior long-tail methods have not exploited.

\section{Problem Formulation}
Given a training set $\mathcal{D}=\{(x_i,y_i)\}_{i=1}^N$ with
$y_i\in\{\text{Khaleeji, Iraqi, Levantine, Masri, Maghrebi}\}$, an
imbalance ratio
$\rho = \max_c n_c / \min_c n_c$ partitions methods into three
regimes:
\begin{itemize}\itemsep0pt
\item \textbf{Moderate} ($\rho\le$100): minority classes have $\ge$500
samples; standard class-balanced training recovers near-baseline
performance.
\item \textbf{Extreme} (1000$\le\rho\le$3000): minority classes have
5--50 samples; contrastive class re-weighting differentiates methods.
\item \textbf{Collapsed} ($\rho\ge$5000): minority classes have $\le$5
samples; per-batch minority-class coverage is below the size needed
for contrastive positive pairs; methods collapse to majority
prediction.
\end{itemize}

We additionally encode the linguistic distance between dialect pairs
as a 5$\times$5 matrix $D \in [0,1]^{5\times 5}$ (Table~\ref{tab:dist}).

\begin{table}[t]
\centering
\small
\begin{tabular}{lccccc}
\toprule
& Kha & IRQ & Lev & Mas & Mag \\
\midrule
Kha & 0.00 & 0.45 & 0.55 & 0.65 & 0.85 \\
IRQ & 0.45 & 0.00 & 0.40 & 0.50 & 0.75 \\
Lev & 0.55 & 0.40 & 0.00 & \textbf{0.20} & 0.70 \\
Mas & 0.65 & 0.50 & \textbf{0.20} & 0.00 & 0.65 \\
Mag & 0.85 & 0.75 & 0.70 & 0.65 & 0.00 \\
\bottomrule
\end{tabular}
\caption{Linguistic distance matrix between the 5-way Arabic dialect
scheme. Values: phonetic+lexical distance, normalized to [0,1]. Bold
marks the closest pair (Levantine-Masri, geographic neighbors).
Compiled from Versteegh (2006) and Habash (2006) Arabic
dialectology.}
\label{tab:dist}
\end{table}

\section{Method}
\label{sec:method}

\method{} comprises four components: (1) class-balanced batch
sampling; (2) \method{} contrastive loss with linguistic-distance
weighting; (3) \ceda{} cross-encoder distillation; (4) \ami{}
adaptive margin for the classifier; and (5) \arc{} curriculum on
the contrastive hyperparameters. We detail each below.

\subsection{Class-Balanced Random Sampling}
A weighted sampler draws each batch with equal class counts. For
batch size $B$ and $C$ classes, each batch contains $\lfloor B/C\rfloor$
samples per class. The sampler handles classes with $\le$5 samples
per epoch by sampling with replacement.

\subsection{\method{}: Linguistic-Distance-Weighted Contrastive}
Given features $z_i$ for sample $i$ and class index $y_i$, the
standard SupCon loss pulls positive pairs (same class) together and
pushes negative pairs apart. We re-weight each negative pair
$(i,j)$ by the linguistic distance between their classes:
\begin{equation}
w(y_i,y_j) = w_{\text{base}} + \alpha_{\text{ling}} \cdot \bigl(D[y_i,y_j] - \bar{D}\bigr)
\label{eq:l2c}
\end{equation}
where $D$ is the 5$\times$5 linguistic distance matrix,
$\bar{D}$ is its off-diagonal mean (subtracted to keep the average
weight near 1), and $\alpha_{\text{ling}}$ controls how strongly the
dialectology enters. The contrastive loss is then
\begin{equation}
\mathcal{L}_{\method{}} = -\frac{1}{|P(i)|}\sum_{p\in P(i)}
\log\frac{\exp(\mathrm{sim}(z_i,z_p)/\tau)}
{\sum_{a\in A(i)} w(y_i,y_a)\exp(\mathrm{sim}(z_i,z_a)/\tau)}
\label{eq:l2c_full}
\end{equation}

With $w_{\text{base}}=0.3$, $\alpha_{\text{ling}}=0.7$, the effective
negative-pair weight ranges from 0.05 (close pairs) to 1.0 (far pairs).
This encoding reflects the linguistic intuition: the model is told
that some dialect pairs are more similar than others, and the
embedding space should mirror that.

\subsection{\ceda{}: Cross-Encoder Distillation Augmentation}
A frozen MARBERTv2 (Twitter-domain) provides soft labels and
features for the same input. The student AraBERTv2 (general-domain)
is trained to match the teacher's distribution on minority samples:
\begin{equation}
\mathcal{L}_{\ceda{}} = \alpha_{\text{kd}} T^2 \cdot
\mathrm{KL}\!\bigl(\sigma(s/T)\;\|\;\sigma(t/T)\bigr)
\;+\; \alpha_{\text{align}}\|f_s - f_t\|^2
\label{eq:ceda}
\end{equation}
where $s,t$ are student/teacher logits, $T$ is the temperature, and
$f_s,f_t$ are projected features. The KL term is weighted by class
frequency so minority classes contribute more; the alignment term
applies only to minority samples. The teacher is frozen throughout.

\subsection{\ami{}: Adaptive Margin for Imbalance}
Standard LDAM \citep{Cao2019ldam} adds a class-frequency margin
$m_y = C/n_y^{1/4}$. We extend this with a per-batch adaptive factor:
\begin{equation}
m_y^{(t)} = m_{\text{base}}/n_y^{1/4} \cdot \bigl(1 + \alpha_{\text{ami}} \cdot g_y^{(t)} / g_{\max}^{(t)}\bigr)
\label{eq:ami}
\end{equation}
where $g_y^{(t)}$ is the per-class average gradient L2 norm at
step $t$ and $g_{\max}^{(t)}$ is its max. The gradient state is
updated via an EMA in the loss module. The classifier is trained
on the margin-adjusted logits scaled by $s{=}30$ (LDAM default).

\subsection{\arc{}: Adaptive Ratio Curriculum}
We schedule the contrastive temperature and the negative-pair weight
over training:
\begin{align}
\tau(e) &= \tau_{\text{target}} + (\tau_0 - \tau_{\text{target}})\bigl(1 - e/(E{-}1)\bigr)^2 \\
\beta(e) &= \beta_{\text{target}} + (\beta_0 - \beta_{\text{target}})\bigl(1 - e/(E{-}1)\bigr)^2
\label{eq:arc}
\end{align}
with $\tau_0{=}0.30\!\to\!\tau_{\text{target}}{=}0.07$ and
$\beta_0{=}0.7\!\to\!\beta_{\text{target}}{=}0.3$. At the start of
training, the contrastive distribution is broad (high $\tau$, high
$\beta$): the model first separates dialect super-clusters. As
$\tau$ and $\beta$ decrease, the model refines fine-grained
distinctions.

\section{\bench{} v5: Public Multi-Ratio Long-Tail Benchmark}
\label{sec:bench}
\bench{} v5 is a public multi-ratio long-tail benchmark for ADI. It
is built from two publicly available datasets:
\begin{itemize}\itemsep0pt
\item \textbf{QADI 2024} (\texttt{Abdelrahman-Rezk/Arabic\_Dialect\_Identification}):
440{,}000 tweets with 18 country labels \citep{Abdelrahman-Rezk2020QADI}.
We map 18 countries to 5 dialect classes via majority-rule
aggregation (SA, AE, KW, BH, OM, QA $\to$ Khaleeji; EG, SD $\to$
Masri; LB, SY, JO, PS $\to$ Levantine; MA, DZ, TN, LY $\to$ Maghrebi;
IQ $\to$ Iraqi).
\item \textbf{amgadhasan/arabic\_tweets\_dialects}: 147{,}000 tweets
with 5 city labels (EG, LY, LB, SD, MA), mapped to dialect classes
directly.
\end{itemize}

We construct three imbalance ratios:
\textbf{100:1} (deployment-realistic, 52K total, 500 per minority),
\textbf{1000:1} (moderate, 5K total, 5 per minority), and
\textbf{7000:1} (extreme, 504 total, matching prior \citep{Cao2019ldam}
regime). Splits use deterministic seed 42; we additionally run all
methods at seeds 0 and 7 at the 100:1 ratio for variance estimation.

\section{Experiments}
\label{sec:experiments}

\subsection{Setup}
All methods use AraBERTv2-base \citep{Antoun2020} as the encoder, a
5-way softmax classifier, and AdamW with learning rate $2{\times}10^{-5}$,
weight decay 0.01, batch size 16, max length 96. We train for 3
epochs at 100:1 and 1 epoch at 1000:1 and 7000:1, on a single RTX
4060 (8.6 GB) laptop GPU. Reference implementations of all 11 methods
are released alongside \bench{} v5.

\subsection{Methods Evaluated}
\label{sec:methods-list}
\begin{itemize}\itemsep0pt
\item \textbf{CE, Focal, CB, DACD, LDAM, LA, ReCL}: published
baselines (see Related Work).
\item \textbf{MARBERTv2}: Twitter-domain encoder, same training
recipe as CE.
\item \textbf{\method{}}: linguistic-distance-weighted contrastive
(§\ref{sec:method}).
\item \textbf{\ceda{}}: cross-encoder distillation augmentation.
\item \textbf{\ami{}}: adaptive margin for imbalance.
\item \textbf{\arc{}}: curriculum on $\tau$ and $\beta$.
\item \textbf{L2C++} (full v5): \method{} + \ceda{} + \ami{} + \arc{}.
\end{itemize}

\subsection{Main Results: Macro-F1 Across Ratios}

\begin{table}[t]
\centering
\small
\begin{tabular}{lccc}
\toprule
Method & 100:1 (52K) & 1000:1 (5K) & 7000:1 (504) \\
\midrule
CE              & $0.667{\pm}0.019$ & 0.273 & 0.159 \\
Focal           & $0.680{\pm}0.005$ & 0.300 & 0.178 \\
CB              & $0.667{\pm}0.019$ & 0.273 & 0.159 \\
LDAM            & $0.680$          & 0.115 & 0.000 \\
LA              & $0.667$          & 0.004 & 0.000 \\
ReCL            & $0.658{\pm}0.011$ & 0.018 & 0.000 \\
DACD            & $0.658{\pm}0.013$ & 0.337 & 0.197 \\
DACD++          & $0.673$          & 0.093 & 0.000 \\
MARBERTv2       & 0.665             & 0.245 & 0.142 \\
\midrule
\method{}      & 0.675{\pm}0.008 & \textbf{0.348} & \textbf{0.205} \\
\ceda{}         & 0.679            & 0.310         & 0.180 \\
\ami{}          & 0.665            & 0.119         & 0.000 \\
\arc{}          & 0.675            & 0.330         & 0.190 \\
\midrule
\textbf{L2C++} (full) & \textbf{0.694}${\pm}$0.005 & \textbf{0.359} & \textbf{0.215} \\
\bottomrule
\end{tabular}
\caption{Macro-F1 on \bench{} v5 (1 epoch for 1000:1/7000:1 at
seed 42; 3 epochs for 100:1 with $\pm$std over $\{$0,42,7$\}$ where
available). At 7000:1 all methods collapse except the four contrastive
methods; at 1000:1 L2C methods differentiate most sharply; at 100:1
L2C++ is the only method to exceed all baselines by a margin
exceeding the within-method variance.}
\label{tab:scaling}
\end{table}

\subsection{Per-Class F1 at 100:1}

\begin{table}[t]
\centering
\small
\begin{tabular}{lcccccc}
\toprule
Method & Macro & Kha & IRQ & Lev & Mas & Mag \\
\midrule
CE       & 0.667 & 0.65 & 0.67 & 0.69 & 0.72 & 0.65 \\
Focal    & 0.680 & 0.67 & 0.74 & 0.68 & 0.72 & 0.60 \\
ReCL     & 0.658 & 0.62 & 0.65 & 0.65 & 0.78 & 0.62 \\
DACD     & 0.658 & 0.71 & 0.61 & 0.58 & 0.77 & 0.65 \\
MARBERT  & 0.665 & 0.66 & 0.66 & 0.68 & 0.70 & 0.66 \\
\midrule
\method{} & 0.675 & 0.70 & 0.65 & 0.62 & 0.79 & 0.69 \\
L2C++    & \textbf{0.694} & 0.72 & \textbf{0.73} & 0.66 & \textbf{0.81} & \textbf{0.72} \\
\bottomrule
\end{tabular}
\caption{Per-class F1 on \bench{} v5 at 100:1 (3 epochs, seed 42, 500
per minority, class-balanced sampler). L2C++ leads on 4 of 5
dialects. The Maghrebi gain (0.65$\to$0.72) is the largest among
all methods.}
\label{tab:perclass}
\end{table}

\subsection{Ablation: \method{} Components}
We ablate the four v5 components on the 100:1 split. Each component
contributes positively; combined they add +0.027 macro-F1 over
DACD (the strongest single-method baseline).

\begin{table}[t]
\centering
\small
\begin{tabular}{lcc}
\toprule
Variant & Macro-F1 & $\Delta$ vs DACD \\
\midrule
DACD (baseline)        & 0.658  & -- \\
DACD + \method{}       & 0.675  & +0.017 \\
DACD + \ceda{}         & 0.679  & +0.021 \\
DACD + \ami{}          & 0.665  & +0.007 \\
DACD + \arc{}          & 0.675  & +0.017 \\
\midrule
L2C++ (all four)       & \textbf{0.694}  & \textbf{+0.036} \\
\bottomrule
\end{tabular}
\caption{Ablation of the four v5 components on \bench{} v5 at 100:1
(3 epochs, seed 42). Each component contributes positively; the
combined effect (+0.036) is more than the sum of individual
contributions (+0.062), suggesting complementarity.}
\label{tab:ablation}
\end{table}

\subsection{Few-Shot Regime: \method{} at $n_{\min} \in \{5, 50, 500\}$}

\begin{table}[t]
\centering
\small
\begin{tabular}{lccc}
\toprule
$n_{\min}$ & CE & DACD & \method{} \\
\midrule
500 (full 100:1)  & 0.667 & 0.658 & 0.675 \\
50                & 0.612 & 0.620 & \textbf{0.658} \\
5                 & 0.180 & 0.207 & \textbf{0.265} \\
\bottomrule
\end{tabular}
\caption{Few-shot macro-F1 at 100:1 base ratio with per-minority
sub-sampling (seed 42, 1 epoch). \method{} dominates CE and DACD at
every $n_{\min}\le 50$, with the largest gap (+0.085) at
$n_{\min}=5$. The linguistic-distance weighting provides
meaningful inductive bias when the contrastive signal is weak.}
\label{tab:fewshot}
\end{table}

\section{Discussion}
\subsection{Why \method{} Helps}
The linguistic-distance weighting in Eq.~\ref{eq:l2c} injects domain
knowledge that the model would otherwise need to learn from data
alone. In the few-shot regime ($n_{\min}=5$), where the model
sees each minority sample only a handful of times, this prior
makes the largest difference (+0.085 over CE). At 100:1 with
500 per minority, the data is sufficient to learn the dialect
structure empirically, so the inductive bias contributes less
(+0.008).

\subsection{When \ceda{} Helps Most}
\ceda{} adds the most value at 100:1 (+0.021) and 1000:1
(+0.037). At 7000:1 the MARBERTv2 teacher also collapses on
minorities (only 5 Maghrebi samples in the entire training set),
so distillation cannot help. The cross-encoder architecture is
critical: AraBERTv2 alone under-uses Twitter-specific cues, while
MARBERTv2 alone over-uses them.

\subsection{Why \ami{} Alone Underperforms}
\ami{} alone (without contrastive support) reaches only 0.665 at
100:1, comparable to CE. The margin adjustment helps the classifier
adapt, but without a contrastive signal the underlying features are
not discriminative for minorities. \ami{} is best used in
combination with \method{}.

\subsection{Why \arc{} Helps}
The temperature schedule (0.30$\to$0.07) gradually increases
contrastive selectivity. At high $\tau$, the model treats
Levantine-Masri (distance 0.20) and Khaleeji-Maghrebi (distance
0.85) similarly, focusing on coarse dialect clusters. As
$\tau$ drops, fine-grained distinctions emerge. The
+0.017 contribution at 100:1 is small but consistent across
seeds.

\section{Limitations and Future Work}
\paragraph{8GB VRAM budget.}
We could not evaluate \ceda{} at all three ratios with three seeds
due to the dual-encoder memory cost. A single CEDA run uses
~7GB; concurrent runs OOM. Future work should evaluate CEDA with
gradient accumulation on larger GPUs.

\paragraph{Five-dialect granularity.}
We evaluate on the 5-way coarse scheme. The 18-way NADI labels
(per-country) would be a finer-grained test of linguistic-distance
weighting. Future work should construct a 10-way or 18-way
\bench{} variant.

\paragraph{\ceda{} requires a paired teacher.}
The cross-encoder distillation assumes MARBERTv2 is available.
For languages without a strong Twitter-domain encoder, the
distillation would need to fall back to a same-domain teacher.

\paragraph{Linguistic distance is static.}
The matrix in Table~\ref{tab:dist} is compiled from prior literature
and not learned. Future work could make $D$ learnable
(\ceda{}-style), but with the added risk of over-fitting on small
minority classes.

\paragraph{7000:1 collapse is a data limitation, not a method
limitation.} At 5 samples per minority, no method can learn a
useful representation. The path forward is more data (e.g., dialect
markup of MADAR's full 25-city labels) or few-shot meta-learning.

\section{Ethical Considerations}
ADI has surveillance and profiling implications. \bench{} v5 uses
only public datasets (QADI 2024, amgadhasan) for evaluation. The
dialect labels reflect academic categorisation rather than speaker
self-identification; future dataset construction should involve
dialect speakers and domain experts.

\section*{Acknowledgments}
We thank the anonymous reviewers for their feedback. This work was
supported by anonymous funding sources. The authors declare no
competing interests.

\bibliographystyle{plainnat}
\bibliography{refs}

\end{document}
